\section{Grammaticality and acceptability}\label{sec:gramm-accept}

The reassessment also needs a disciplined vocabulary. The background view
adopted here is that grammar is a system of projectibility profiles
\citep{Goodman1955}: relatively stable patterns that license explicitly
scoped inferences from forms, meanings, contexts, speakers, and judgments to
uses a community or register treats as available. Research purposes select
which inference is worth testing; they don't make the inference true.

The empirical projective claim in this paper is deliberately narrow. For the
sampled contrasts in the 2013 and 2017 comparison data, observing relative
ordering under one formal elicitation format warrants expecting similar
aggregate ordering under another,
within the variation reported here. A grammaticality claim has wider targets:
interpretation, production, correction, expert classification, register fit,
cross-speaker stability, and related consequences. Those targets require data
beyond the task-format comparisons analysed in this paper.

The paper distinguishes three observation channels from the target. Participant
responses under an acceptability prompt are observations of judgments under a task. Expert
asterisks and published example labels are classifications made within a
theoretical or editorial practice. Benchmark labels are dataset-construction
decisions. The projectibility profile is the target: the structure of inferences
a form licenses, what its acceptability, production, correction, expert
classification, register fit, and cross-speaker stability jointly project to. Two
channels are evidence about the same profile when they bear on the same form's
inferential relations.

So the profile isn't the average of the channels, and channel disagreement is
not noise to smooth away. When a sentence is rated low by naive participants,
left unstarred in the literature, and attested in a register corpus, the verdict
about its profile is fixed by which inferences actually hold, tested against
further data, not by a vote across channels.

A channel that consistently fails to predict the others establishes divergence,
not its interpretation. Independent evidence has to decide whether the divergence
reflects task specificity, noise, or a distinct part of the profile. That
requirement makes the profile more than an unweighted set of labels and supplies the identity
condition the level discipline needs.

This discipline keeps neighbouring constructs apart. Acceptability is one
evidential route into a projectibility profile, not the profile itself. Personal
production, naturalness, standardness, correction behaviour, and theoretical
classification project differently. A sentence can be dispreferred without
being ungrammatical, correctable without being unnatural, or standard in one
register while unacceptable under another task framing.

The framing earns its place in the measurement argument by predicting where a
single task-conditioned acceptability scalar could be insufficient: items where
participants rate a sentence low but the literature leaves it unstarred, or
where the form is standard in one register and degraded under another task
framing. A dataset that samples those channels jointly should put a
single-scalar model at risk.

The present comparison matrices contain formal task responses rather than that
cross-channel evidence. They test the narrower question of whether ME
carries a stable method-specific dimension or practical ordering that the
bounded tasks miss. The projectibility account supplies a wider empirical
programme, but the paper doesn't treat the task-format result as having already
tested it.

Bard et al.\ helped make one part of this space measurable: graded
acceptability judgments from participants \citep{BardEtAl1996}. Their work made a
major contribution to the measurement of one evidential channel. It doesn't by
itself show that an acceptability scale captures every inference licensed by a
grammar. Later benchmarks sometimes repackage acceptability-related labels as
grammaticality targets \citep{WarstadtEtAl2019CoLA,WarstadtEtAl2020}. The
Syntactic Acceptability Dataset is valuable here because it explicitly
separates literature-derived grammatical status from crowd acceptability status
\citep{juzek2024}.

That boundary also limits the no-new-data paper. Existing datasets can compare
the stability of elicitation methods and clarify how acceptability labels
travel. They can't settle a full projectibility profile for a construction:
production, correction, register, interpretation, and community licensing would
require data designed for those inferences.
